\chapter{Materiales e Instrumental}

Para la ejecución de las actividades de relevamiento de campo, mediciones de laboratorio y ensayos de susceptibilidad
electromagnética, se emplearon los siguientes instrumentos, componentes y dispositivos de prueba:

\section{Instrumental de Medición y Ensayo}
\begin{description}
    \item[Osciloscopio Digital Tektronix TDS 220:] De dos canales, ancho de banda analógico de $\qty{100}{\mega\hertz}$
    y tasa de muestreo de $\qty{1}{\text{GSa/s}}$, utilizado en la cuantificación de las perturbaciones inducidas por
    la red eléctrica de $\qty{50}{\hertz}$.
    \item[Osciloscopio Digital RIGOL DS1052E:] De dos canales, ancho de banda de $\qty{50}{\mega\hertz}$, tasa de
    muestreo de $\qty{1}{\text{GSa/s}}$, utilizado en los ensayos del banco de celda de carga.
    \item[Puntas de Prueba:] Puntas pasivas compensadas con atenuación conmutable 1X/10X y terminal con cable de tierra
    de bajas dimensiones inductivas.
    \item[Multímetro Digital:] Multímetro digital para la determinación la tensión de salida continua de los sensores.
    \item[Fuente de Alimentación de CC:] Fuente de laboratorio regulada lineal, ajustada a $\qty{+9}{\volt}$ para
    alimentar la placa de acondicionamiento.
    \item[Generador de EMI:] Herramienta rotativa tipo torno de mano (Dremel) con motor universal de colector a delgas
    y escobillas alimentado desde la red de $\qty{220}{\volt}$, generador de ruido electromagnético de alta frecuencia.
\end{description}

\section{Elementos del Banco de Ensayos de Interferencia}
\begin{description}
    \item[Celda de Carga:] Celda de carga de aluminio tipo viga a flexión de $\qty{3}{\kilogram}$ de fondo de escala,
    puente de Wheatstone completo con galgas extensométricas de alta estabilidad y cable blindado con malla metálica.
    \item[Placa de Acondicionamiento:] Placa de circuito impreso basada en amplificador de instrumentación
    de precisión (INA122/INA128), regulador de tensión local MC78L05 ($\qty{+5}{\volt}$) para excitación
    ratiométrica del puente, resistencia de ganancia $R_G$ y capacitores cerámicos de desacoplo de
    $\qty{100}{\nano\farad}$ y $\qty{10}{\micro\farad}$.
    \item[Gabinete Metálico:] Chasis metálico conductor que actúa como blindaje electrostático (jaula de Faraday), con
    terminal accesible para derivación de corrientes parásitas hacia el potencial de referencia (GND).
    \item[Conector Auxiliar DB15:] Conector macho cableado para interconectar directamente los bornes diferenciales de
    la celda de carga al osciloscopio y multímetro sin pasar por la etapa de ganancia.
    \item[Cargas Patrón de Ensayo:] Batería sellada y juego de destornilladores en caja plástica.
\end{description}

\section{Sistemas Relevados en Planta}
\begin{description}
    \item[Máquina de Tracción y Compresión:] Sistema hidráulico industrial del Laboratorio de Ensayos del
    Departamento de Ingeniería Mecánica de la UTN-FRC, equipado con transductor de desplazamiento lineal
    potenciométrico GEFRAN modelo PY-2-C-100 ($\qty{5}{\kilo\ohm}$, carrera de $\qty{100}{\milli\meter}$,
    linealidad independiente $\pm 0.07\%$) y transmisor de presión ADZ NAGANO Sensortechnik
    ($\qtyrange{0}{250}{\bar}$, salida $\qtyrange{0}{10}{\volt}$).
    \item[Banco de Ensayo de Motores:] Banco dinamométrico del Laboratorio de Ensayos de Motores de la UTN-FRC,
    compuesto por freno dinamométrico oscilante con brazo de palanca instrumentado mediante celda de carga de
    precisión tipo S-beam (puente de Wheatstone completo de galgas de $\qty{350}{\ohm}$, sensibilidad
    $\qty{3}{\milli\volt\per\volt}$) y sensor de efecto Hall enfrentado a rueda fónica dentada (encoder de RPM).
\end{description}
